% Part: many-valued-logic
% Chapter: syntax-and-semantics
% Section: formulas

\documentclass[../../../include/open-logic-section]{subfiles}

\begin{document}

\olfileid{mvl}{syn}{fml}

\olsection{\usetoken{P}{formula}}

\begin{defn}[సూత్రం]
\ollabel{defn:formulas}
ప్రతిజ్ఞావాక్య భాష~$\Lang L$లోని \emph{\tetoken{సూత్రాల}{!!{formula}s}}
సమితి~$\Frm[L]$ను కింది విధంగా ఆగమనాత్మకంగా నిర్వచిస్తాం:
\begin{enumerate}
\item ప్రతి \tetoken{ప్రతిజ్ఞావాక్య చరం}{!!{propositional variable}}~$\Obj p_i$
  ఒక పరమాణు \tetoken{సూత్రం}{!!{formula}}.
\item $\Lang L$లోని ప్రతి $0$-స్థానిక సంయోజకం (ప్రతిజ్ఞావాక్య స్థిరాంకం)
ఒక పరమాణు \tetoken{సూత్రం}{!!{formula}}.
\item $\star$ అనేది $\Lang L$లోని $n$-స్థానిక సంయోజకమై, $!A_1$,
\dots, $!A_n$లు \tetoken{సూత్రాలు}{!!{formula}s} అయితే,
$\star(!A_1, \dots, !A_n)$ కూడా \tetoken{ఒక సూత్రం}{!!a{formula}}.
\tagitem{limitClause}{ఇవి తప్ప మరేదీ \tetoken{సూత్రం}{!!a{formula}} కాదు.}{}
\end{enumerate}
$\star$ $1$-స్థానికమైతే, $\star(!A_1)$ను తరచుగా
$\star !A_1$ అని సంక్షిప్తంగా రాస్తాం. $\star$ $2$-స్థానికమైతే,
$\star(!A_1,!A_2)$ను తరచుగా $(!A_1 \star !A_2)$గా రాస్తాం.
\end{defn}

మామూలుగానే, బయటివైపు ఉన్న కుండలీకరణలను చాలాసార్లు చెప్పకుండానే వదిలేస్తాం.

\begin{ex}
  ప్రామాణిక భాష~$\Lang{L_0}$లో, $\Obj p_1 \lif (\Obj p_1
  \land \lnot \Obj p_2)$ ఒక సూత్రం. గుణిత తర్కపు భాషలో దాన్ని
  $\Obj p_1 \lif (\Obj p_1 \odot
  \lnot \Obj p_2)$గా రాస్తాం. భాషకు $1$-స్థానిక
  $\triangle$ను చేరిస్తే,
  $\triangle (\Obj p_1
  \land \Obj p_2) \lif (\triangle \Obj p_1 \land \triangle \Obj p_2)$
  వంటి సూత్రాలూ ఉంటాయి.
\end{ex}

\end{document}
